\chapter{Exceptions}

% normal flow: the normal flow of a program is that one statement is evaluated after another, some are conditionally evaluated and some are repeated, functions are called (parameters transferred to them, work on the stack to keep track of local caller/callee variables), these functions return the thread of execution to the caller that then carries on with its evaluation, eventually the last statement will complete and the program will exit
The normal flow of a program is that one statement is evaluated after another; some are conditionally evaluated, some are repeated, and functions are called. Parameters are transferred to the invoked function and work is being done one \idx{the stack}{Stack!The} to keep track of local \idx{caller}{Caller} and \idx{callee}{Callee} variables. Once the end of the function is reached, the \idx{thread of execution}{Thread of execution} is returned to the caller that then carries on with its evaluation. Eventually, the last statement will complete and the \idx{process}{Process} (or program) will \idx{exit}{Exit}.

% error reporting through return values: certain operations cannot be guaranteed to be successful, functions implementing such operations will have some mechanism for reporting whether they were successful, historically this has often been done through the return value and/or global variables, global variables have to downside of being shared across all threads of execution (and that can make it hard to interpret), this is not a concept that we cover in this book (but it is highly relevant soon thereafter), returning an error state is very common but most programming languages only supports returning a single value, so one would have to find a concrete value to represent an error, if you want to read a number of bytes from a file (that doesn't exist) and the function should return the number of bytes actually read then any successful operation will result a non-negative number of bytes read (and that means that we can use negative values to represent errors), this is cumbersome though and we have to keep the error coding present in our minds, it also forces us to check the error after each function call (that could produce an error), the handling is coupled to the function call
Certain operations cannot be guaranteed to be \idx{successful}{Success}. Functions implementing such operations will have some \idx{mechanism}{Mechanism} for reporting whether they were successful. Historically, this has often been done through the \idx{return value}{Return value} and/or global variables. The choice of \idx{global variables}{Variable!Global} has the downside of them being shared across all \idx{thread of execution}{Thread of execution} (and that make it hard to interpret). \idx{Threading}{Threading} is not a concept that we cover in this book, but it will likely become highly relevant soon after you finish it. Returning an error state is very common, but as most \idx{programming languages}{Programming language} only support returning a single value, one would have to find a concrete value to represent an error. If you want to read a number of bytes from a file (that doesn't exist) and the function should return the number of bytes actually read, then any successful operation will return a non-negative number representing the number of bytes read. This means that we can use negative values to represent \idx{errors}{Error}. This is cumbersome though, and we have to keep the error \idx{coding}{Error!Coding} present in our minds. It also forces us to \idx{check}{Error!Checking} the error after each function call (that could produce an error). The handling of the error is \idx{coupled}{Coupling} to the function call.

% implicit partial handling with exceptions: some languages support the notion of an exception, exceptions decouple the discovery that something has gone wrong from the handling and provides a means for transporting the error code (and additional information) from the site of discovery to the site of handling, with exceptions we no longer have to encode error indicators in return values or global variables
Some languages have support for \idx{exceptions}{Exception}. An exception is a different mechanism for handling such abnormal flows. They decouple the discovery that something has gone wrong from the handling of that situation. They also provide a means for transporting the error code (and additional information) from the \defi{site of discovery}{Site!Exception discovery} to the \idx{site of handling}{Site!Exception handling}. With exceptions, we no longer have to encode error indicators in return values or global variables.

\section{Lifecycle of an Exception}

% the stages of the life of an exception: construction, throwing, handling
In order to discuss what an exception is, we will take a look at its lifecycle. As shown in figure \ref{fig:exception:lifecycle}, the life of an exception goes through three phases, namely:
\begin{enumerate}
  \descitem{Construction:} An exception is a form of data structure, and it has a type. We first construct an exception, and then we can do operations on it.
  \descitem{Throw:} One of those operations is to \textsl{throw} it. This is what we do when we discover that something has gone wrong, and it will abruptly end the normal flow of execution. We will refer to this as the \defi{production site}{Site!Exception production} of the exception.
  \descitem{Catch:} At an earlier point on our execution thread we have defined a \defi{handling site}{Site!Exception handling} (aka \defi{catch site}{Site!Exception catch}). This is kind-of a checkpoint. The latest defined handling site that matches the thrown exception will \textsl{catch} it. That means that execution will continue at this location.
\end{enumerate}
\textbf{Note:} The handling site has access to the exception data structure and has the option of throwing it once again.

% figure: looping throwing and handling
\begin{figure}[tbp]
  \begin{center}
  \begin{tikzpicture}[]
    \newcommand{\spacing}[0]{12mm}
    
    \tikzstyle{edge}  = [thick,>=stealth,draw=black]
    \tikzstyle{dedge} = [thick,->,>=stealth,draw=black]
    \tikzstyle{node}=[
      rectangle,
      draw=purple,
      anchor=west,
      thick,
    ]
    
    \node[node] (construction) at (0, 0) {Construction};
    \node[node] (throw) at ([xshift=\spacing]construction.east) {Throw};
    \node[node] (catch) at ([xshift=\spacing]throw.east) {Catch};
    
    \draw[dedge] (construction)--(throw);
    \draw[dedge] (throw)--(catch);
    \draw[dedge,dotted] (catch)
              -| ([xshift=\spacing,yshift=\spacing/2]catch.north east)
              -| (throw);
  \end{tikzpicture}
\end{center}

  \caption{Lifecycle of an exception.}
  \label{fig:exception:lifecycle}
\end{figure}

% locating the site of handling: imagine that throughout the execution of our code we keep track of all the handling sites that we come across, whenever an exception is thrown it is caught by the code of the latest handling site, execution immediately jumps to this point in the code (even if it is not part of the current function invocation, in that case the environment of that function is reestablished from the stack)
We can have multiple handling sites for the same exception. So, which one is used? To answer this, let's imagine that throughout the execution of our code we keep track of all the handling sites that we come across. Whenever an exception is thrown, it is caught by the handler (i.e., code) of the latest handling site. \idx{Execution}{Execution} immediately jumps to this point in the code. If it is not part of the current \idx{function invocation}{Function!Invocation}, \idx{stack frames}{Stack frame} will be popped off until the one matching the handling site is reached, and the \idx{environment}{Environment} stored here will be reestablished.

\section{\csharp}

% intro: how to go about it
To cover exceptions in \csharp, we need to introduce syntax \ref{syntax:exception}. That is a lot of rules! However, when we group them by use, we can focus on a one or two of them at a time. So, let's do that.

\begin{syntaxfloat}
  \input{syntax/exception.tex}
  \caption{Exceptions.}
  \label{syntax:exception}
\end{syntaxfloat}

\subsection{Creating Exceptions}

% new: what is the new keyword used for?

\inputminted[]{csharp}{../src/csharp/exceptions/creation/Program.cs}

\subsection{Throwing Exceptions}

% intro: once we have an exception we can throw it, for now we use the expr track of the stmt rule from syntax \ref{syntax:exception}. It looks like this:
Once we have an exception, we can throw it. For now, we use the \conceptname{expr} track of the \conceptname{stmt} rule from syntax \ref{syntax:exception}. It looks like this:

\inputminted[firstline=1]{csharp}{../src/csharp/exceptions/throwing/Program.cs}

% what happens
At the moment we have no code to catch the exception. So, when the exception occurs, the program crashes. This causes the exception to be printed:

\begin{minted}{shell-session}
$ dotnet run
Unhandled exception. System.Exception: Oops!
   at Program.<Main>$(String[] args) in Program.cs:line 2
\end{minted}

% how to interpret
A few things are going on here. Let's go through them one by one:
\begin{itemize}
  \item The virtual machine will print all unhandled exceptions. An unhandled exception, is an exception that have not been caught by handling site and has reached the virtual machine.
  \item The exception has the type \typename{System.Exception}, for which \typename{Exception} is a a shorthand.
  \item Attached to the exception was a \valuename{"Oops!"} payload.
  \item The last line is the stack trace. It provides us with a path from the button of the stack to the production site. For each step we get the function (\keywordname{Main}), the file (\filename{Program.cs}) and the line in this file (line \valuename{2}). In a larger program, the stack trace will typically take up more lines; one for each stack frame.
\end{itemize}

\subsection{Catching Exceptions}

% intro: any uncaught exception will make the program crash, at times this might be what we want  but often we are trying out some approach and then want to react if an exception is thrown, for catching exceptions we rely on the stmt rule of syntax \ref{syntax:exception}, the following example is close to the simplest possible illustration of the mechanism:

\inputminted[firstline=1]{csharp}{../src/csharp/exceptions/catching/Program.cs}

% result: executing that code will result in the following. Please note that the last printout makes it to the output so it must have been evaluated, so catching the exception stops it from propagating

\begin{minted}{shell-session}
$ dotnet run
System.Exception: Oops!
   at Program.<Main>$(String[] args) in Program.cs:line 2
But still alive ...
\end{minted}

% exception can be produced further up the stack

\inputminted[firstline=1]{csharp}{../src/csharp/exceptions/catching_fun/Program.cs}

% result: when we execute this code we see mostly the same behavior, but since the production site is one stack frame further up we get one extra line in the stack trace:

\begin{minted}{shell-session}
$ dotnet run
System.Exception: Oops!
   at Program.<<Main>$>g__fun|0_0() in Program.cs:line 2
   at Program.<Main>$(String[] args) in Program.cs:line 6
But still alive ...
\end{minted}

\subsection{Rethrowing}

% intro: sometimes we may want to know whether an excepton has happened without having to deal with it, notice how when we caught an exception we got a variable of type \typename{Exception} and how that was exactly what we threw? Well, we can throw it again, so we catch the exception react to it and throw it again, the code looks a bit different that you would assume:
At times we may want to know whether an exception has occurred, but we don't want to have to deal with it. To do this, we first notice how when we caught an exception we got a variable of the same \idx{type}{Type} as the expression we threw. That means that we can throw it again! So, we catch the exception, react to it, and then throw it again. When we catch an exception and throw it again, we say that we \idx{rethrow}{Exception!Rethrowing} it. The practical code looks a bit different than what this suggests:

\inputminted[firstline=1]{csharp}{../src/csharp/exceptions/rethrowing/Program.cs}

% note: note how throw is used without an explicit reference to the exception that we want thrown? When throwing an exception the stack trace part of that exception is updated to reflect the location in the code where the exception is thrown, we typically recast to have a peak at what is going on, updating the stack trace would make it look like the exception originated from the part of the code that is peaking and that would make it difficult to debug, we would know what has happened but not where, \csharp\ has this implicit throw mechanism to rethrow the currently caught exception without updating its stack trace, accordingly we don't need to bind a name to that matched exception, in fact we would get a warning if we tried, this because if would most likely lead to trouble
Note how \keywordname{throw} is used without an explicit reference to the exception that we want to throw? When throwing an exception, the \idx{stack trace}{Stack trace} field of that exception is updated to reflect the location in the code where the exception is thrown. We typically rethrow to have a peak at what is going on. Updating the stack trace would make it look like the exception originated from the part of the code that is peaking, and that would make it difficult to \idx{debug}{Debug}: While we would know what happened, we would not know \textsl{where} it happened. \csharp\ has this implicit variant of the \keywordname{throw} mechanism to rethrow the last caught exception without updating its stack trace. Accordingly, we don't need to \idx{bind}{Variable!Binding} a name to that matched exception. In fact, as it would most likely be a \idx{typo}{Typo}, we would get a warning if we tried. Why else would we \idx{declare}{Variable!Declaration} a variable without using it?

% explanation
When we execute the program we see that the code in our handler is evaluated, and that the code still crashes as if we hadn't caught the exception:
\begin{minted}{shell-session}
$ dotnet run
Things are about to go badly ...
Unhandled exception. System.Exception: Oops!
   at Program.<Main>$(String[] args) in Program.cs:line 2
\end{minted}

\subsection{Finalization}

\inputminted[firstline=1]{csharp}{../src/csharp/exceptions/finalization/Program.cs}

\begin{minted}{shell-session}
$ dotnet run
System.Exception: Oops!
   at Program.<Main>$(String[] args) in Program.cs:line 3
Anyways ...
1
Anyways ...
\end{minted}

\section{C}

\section{Go}

\exercises{exceptions}{Exceptions}
